
\section{Introduction}
\label{sec:intro}

%\begin{figure}[t]
%  \begin{center}
%    \includegraphics[width=0.95\columnwidth]{figs/workflow.eps}
%    \caption{Workflow}
%    \label{fig:Normal workflow for most multi-tier service oriented systems}
%  \end{center}
%\end{figure}


As application software grows and gets more complicated, testing large scale applications has become increasingly important. 
However, it is often impossible to recreate realistic workloads in an offline development environment for large scale multi-tier or cloud based applications.
Testing in the development environment can be (1). Un-realistic because it may not be possible to faithfully reconstruct the production environment, (2). The test-cases generated may be incomplete, (3) It is infeasible to test all possible configurations given time and cost constraints of releasing the software to the field. 

One of the proposed mechanisms of addressing this problem is to ``perpetually test''\cite{perpetual} the application in the field after it has been deployed. 
This is important since testing in a production system enables us to capture previously ``unreachable'' system states, which are possible only for a long running production environment.
%Since we are testing in a real environment, we get real user-input for test-cases, and are able to test a real environment.
%However, perpetual testing has never gotten much traction because of an obvious flaw : executing such test-cases will adversely effect the user-experience, both in performance and potentially in application logic.
Companies such as Flickr, Twitter, Facebook and Google have re-ignited the debate towards similar approaches as perpetual testing under the umbrella of DevOps.
DevOps stresses on close coupling between software developers and operators, and to merge the operations of both. 
Most of these companies have very frequent releases and a very short time to a bug fix, test, patch and release in order to realize continuous delivery(Facebook mobile has 2 releases a day, and Flickr has 10 deployment cycles per day).
Testing, analyzing and monitoring are an important aspect of the devops cycle, however they have several limitations as only a very small performance bandwidth can be dedicated to QA operations as compared to real user activity(so as to not effect user-perceived delay).

%Embedding test-case logic within the context of the application will result in state-change, or performance slow-down that would otherwise not be there in an optimized implmentation of the application.
%This is something which is usually unacceptable in user-facing applications.
%The authors have previously looked into amortizing the cost of running the test, by running it in parallel \cite{invite}.
%However such approaches cannot completely avoid the slowdown, and additonally do not completely sandbox the effects of the test case on the production server. 

%An alternate approach is to record live applications, and replay them offline.
%Over the years there have been several systems which have explored this direction with promising results.
%However most record and replay systems have high overheads, and require replication of the running configuration which may not be possible. 
%Additionally the administrator needs to wait for offline analysis, instead of doing real-time diagnosis.

The motivation behind our work is to provide testing as a service for real-time diagnosis of production applications, in order to signifcantly reduce the time towards bug resolution.
We observe that most modern day service oriented applications are hosted on IAAS cloud providers, and can hence be easily scaled  up. 
%On the other hand, there has been an impressive increase in the scale of computing resources, and distributed scalability of infrastructure.
%Web based applications are often hosted in cloud environments, this allows for easily scaling up the hardware resources.
%This often allows for redundant computation, which can be used for testing purposes. 
Leveraging this abundance of resource, and recent advances in user-space container technology(OpenVZ/LXC\cite{openvz,linux}) we present a testing harness which allows the user to dynamically insert test cases in a production environment(we call this in-vivo testing), enabling real-time diagnosis.
Our system called \texttt{Parikshan}\footnote{Parikshan is the sanskrit word for testing} allows capturing the context of application, and for tests to be run without effecting the sanctity and performance of the actual user-facing application. 
This is done by cloning a production server and creating two containers: a production container, and a testing container. 
We duplicate the incoming traffic to both the production container and the test container using a custom proxy, which ignores the responses from the test-container. 
The testing on the test-container is done on the fly using dynamic instrumentation, hence any set of test-cases can be turned on whenever required. 
%This is achieved by using dynamic instrumentation mechanisms to clone a VM by forking off from a running executed state and encapsulating the forked execution in a VM.
The user can pre-define probe points for dynamically inserting test-cases (by default the entry and exit of each function is considered a probe point).
Since the test is executed in a VM it acts like a sandbox which restricts it from causing any perturbation to the state of the parent process, or effecting the sanity of the responses to the production client. 
We synchronize the production and test containers using a variant of live migration without suspending the services of the production server, and follow it up with frequent synchronization for a long running tests. 
%Frequent synchronization is necessary because the test container can potentially go out of sync with the production because of non-determinism or because a test-case changes the state of the container, and effects future problems. 

\textit{@Nipun edit -> consider why use user-space containers instead of VMs?}
While VM virtualization has existed for several years, recent advances in user-space container technologies, along with support for migration, has created a space for light-weight testing in live environments.
Technically our sandbox techniques could also be applied using more traditional Virtual Machines. 
However, the overhead of using Virtual Machines is considerably higher, and it would technically require double the amount of resources for the target production servers.
User-Space containers reduce this overhead considerably by using the resources in the same machine.
We believe the availablity of resources in IAAS cloud infrastructures combined 


The key contributions of this paper are:

\begin{itemize}
\item A tool which provides a sandbox environment to execute test cases in the production environment. 
This allows for a safe and secure test harness which does not effect the production state, and allows the application to proceed in it's execution.
\item We allow for dynamic insertion of the test case, and safetly capturing the context of the application. Dynamically inserting test-cases is important to avoid relaunching binaries in the test-container with the required test-cases. 
Restarting binaries is not possible, because it would break active network connections, and destroy the state of the test container
\item Language and Platform agnostic: One of the key advantages of our approach is that it is language and platform agnostic. Since the underlying mechanism takes advantage of containers as a platform to do the cloning, the language or interface does not matter as far as cloning is concerned. 
Of-course testing mechanims may differ depending upon different languages.
%\texttt{Zero-Probe Effect} probe points are added to the application which can be activated to insert test cases using ptrace\cite{ptrace}.
%The use of dynamic instrumentation capability to add test cases in an application is an extension of our previous work of a dynamic instrumentation tool iProbe \cite{iProbe}
\end{itemize}

%Traditional testing approaches break states and are unable to  
%The authors previous work in in-vivo testing\cite{invite} explored testing in the wild by initiating test cases in the production environment and sharing the load across several instances of deployed application.
%This approach adds test-cases in predetermined functions before starting the execution of the process, and periodically executes them in the run-time environment based on a probabilistic function. 

%\cite{dapper}

\subsection{Impact}
\label{sec:impact}

The impact of sandbox testing can be seen in several different ways

\begin{itemize}
  \item \textbf{Sandbox Live Testing}
    One of the key motivations leading to \emph{Parikshan} is to provide a harness to allow the user to test real, live implementations. 


  \item \textbf{Fault Tolerance}
  \item \textbf{Verification}
  \item \textbf{Integration Testing}
\end{itemize}

